\begin{description}
	\item[1行目] ベクトル$\bm{a}=(1,2,1)$を作る
	\item[2行目] ベクトル$\bm{b}=(3,4,2)$を作る
	\item[3行目] 掛け算記号\verb|*|を使っているが，これはベクトルの掛け算ではなく要素ごとの掛け算を意味する。
	\item[4行目] ベクトルの掛け算記号\verb|%*%|を使っており，$\bm{ab}$の計算をしている。この時，ベクトル$\bm{a}$は行ベクトル，$\bm{b}$は列ベクトルと解釈されます。行列計算に適した形であり，デフォルトでは行方向が優位です。
	\item[5行目] ベクトルの掛け算だが\texttt{t()}で転置，すなわち$\bm{b}'$を行っており，行列に適した形に変換されて$\begin{pmatrix}1\\2\\1\end{pmatrix}\begin{pmatrix}3&4&2\end{pmatrix}$と解釈されている。$3\times1$と$1\times3$の計算なのでできあがる行列は$3\times3$のサイズになる。
	\item[6行目] 行列$\bm{A}$を作る
	\item[7行目] 行列$\bm{A}$に対してベクトルをかけているように見えるが，ベクトルの掛け算でなく要素の掛け算記号\verb|*|なので，各行を1倍，2倍，1倍する結果になっている。
	\item[8行目] 行列とベクトルの積$\bm{Aa}$であり，行列計算可能な$3\times3$と$3\times1$と解釈され，結果のサイズは$3 \times 1$になっている。
	\item[9行目] ベクトルと行列の積$\bm{aA}$であり，行列計算可能な$1 \times 3$と$3 \times 3$と解釈され，結果のサイズは$1 \times 3$になっている。
	\item[10-11行目] 行列$\bm{B,C}$を作る
	\item[12行目] 行列の積$\bm{BC}$を計算している。$3\times2$と$2\times3$の行列の積なので，結果のサイズは $3\times3$になっている。
	\item[13行目] 行列の積$\bm{BC'}$を計算しようとしているが，$3\times2$と$3\times2$の積は計算できないので，エラーが返ってくる。
\end{description}
